\chapter{Tổng kết và Hướng phát triển Tương lai}

\section{Kết quả đạt được}
Đồ án đã thiết lập thành công ứng dụng \textbf{Vietnam Climate Pulse}, giải quyết trọn vẹn bài toán trực quan hóa khí hậu Việt Nam giai đoạn 2020-2025 thông qua ba kết quả then chốt. Về mặt dữ liệu, dự án đã tự động hóa thành công quy trình thu thập dữ liệu từ dịch vụ API thời tiết, tiến hành làm sạch và lưu trữ cục bộ dưới định dạng tối ưu để tăng tốc độ truy xuất thông tin. Về trực quan hóa, nhóm đã thiết kế giao diện dashboard theo phong cách ``Editorial Climate Almanac'' tông sáng, nền giấy ấm với hệ màu đất trang nhã, biểu diễn bản đồ địa lý các trạm đo bằng biểu đồ tương tác, kết hợp biểu đồ ma trận phân tích nhiệt độ theo mùa vụ và các biểu đồ tương quan đa biến trực quan sinh động. Về tích hợp trí tuệ nhân tạo, module trợ lý phân tích đã được triển khai thành công, hỗ trợ dịch ngôn ngữ tự nhiên sang câu lệnh truy xuất dữ liệu chỉ đọc để chạy trực tiếp trên hệ thống cục bộ dưới sự kiểm soát và phê duyệt an toàn của con người.

\section{Hạn chế của đồ án}
Mặc dù đạt được đầy đủ các mục tiêu đề ra của môn học trực quan hóa dữ liệu, đồ án vẫn tồn tại một số hạn chế nhất định cần được hoàn thiện. Trước hết, do nguồn dữ liệu thời tiết lịch sử là dạng dữ liệu mô phỏng lưới toàn cầu, hệ thống có thể xuất hiện độ lệch nhỏ so với số liệu thực tế đo đạc trực tiếp bằng nhiệt kế chuyên dụng tại các trạm khí tượng quốc gia của Việt Nam. Thêm vào đó, ứng dụng hiện tại tập trung hoạt động trên môi trường máy tính cá nhân phục vụ mục đích trình diễn và kiểm thử nhanh, do đó chưa tích hợp hệ quản trị cơ sở dữ liệu máy chủ đám mây để phục vụ cho lượng người dùng lớn truy cập đồng thời trong thực tế.

\section{Hướng phát triển tương lai}
Để phát triển đồ án thành một sản phẩm nghiên cứu thực tiễn sâu rộng hơn hoặc ứng dụng thương mại, nhóm đề xuất ba hướng cải tiến trọng tâm. Hướng đi thứ nhất là tích hợp dữ liệu trạm đo thực tế bằng cách liên kết trực tiếp với hệ thống dữ liệu trực tuyến của cơ quan khí tượng thủy văn quốc gia hoặc các trạm đo tự động tại địa phương nhằm bảo đảm độ tin cậy tuyệt đối của số liệu thời gian thực. Hướng đi thứ hai tập trung mở rộng khả năng của mô hình trí tuệ nhân tạo, nâng cấp trợ lý ảo từ dịch ngôn ngữ tự nhiên sang hỗ trợ tự động thiết kế biểu đồ trực quan hoặc kết xuất báo cáo tóm tắt tri thức bằng văn bản tiếng Việt. Hướng đi cuối cùng là triển khai ứng dụng lên môi trường máy chủ đám mây trực tuyến, thiết kế lại lớp cơ sở dữ liệu phân tán kết hợp với các động cơ truy vấn chạy trực tiếp trên trình duyệt để phục vụ rộng rãi cho công chúng.
